\documentclass[xetex,mathserif,serif,handout]{beamer}

\usepackage{xunicode}
\usepackage{xltxtra}
\usepackage{color}
\usepackage{url}
\usepackage{listings}
\usepackage{fontspec}
\usepackage{geometry}
\usepackage{lastpage}
\usepackage{fancyhdr}
\usepackage{amsmath}
\usepackage{amsthm}
\usepackage{amssymb}
\usepackage{blkarray}
\usepackage{multicol}
\usepackage{relsize}

\definecolor{solarized@base03}{HTML}{002B36}
\definecolor{solarized@base02}{HTML}{073642}
\definecolor{solarized@base01}{HTML}{586e75}
\definecolor{solarized@base00}{HTML}{657b83}
\definecolor{solarized@base0}{HTML}{839496}
\definecolor{solarized@base1}{HTML}{93a1a1}
\definecolor{solarized@base2}{HTML}{EEE8D5}
\definecolor{solarized@base3}{HTML}{FDF6E3}
\definecolor{solarized@yellow}{HTML}{B58900}
\definecolor{solarized@orange}{HTML}{CB4B16}
\definecolor{solarized@red}{HTML}{DC322F}
\definecolor{solarized@magenta}{HTML}{D33682}
\definecolor{solarized@violet}{HTML}{6C71C4}
\definecolor{solarized@blue}{HTML}{268BD2}
\definecolor{solarized@cyan}{HTML}{2AA198}
\definecolor{solarized@green}{HTML}{859900}
\definecolor{yaleblue}{HTML}{0E4C92}

\setbeamertemplate{navigation symbols}{}
% \setbeamerfont{title}{family=\old}
% \setbeamerfont{author}{family=\tfont}%
% \setbeamerfont{frametitle}{family=\oldA}
% \setbeamerfont{date}{family=\dfont}

\setbeamertemplate{itemize items}{--}
\setbeamercolor*{item}{fg=black}

\defaultfontfeatures{Mapping=tex-text}
\hypersetup{pdfstartview={FitH}}

\newcommand{\old}[1]{\fontspec[Alternate=1,Ligatures={Common}]{Hoefler Text}\fontsize{18pt}{30pt}\selectfont #1}%
\newcommand{\oldA}[1]{\fontspec[Alternate=1,Ligatures={Common, Rare}]{Hoefler Text}\fontsize{12pt}{15pt}\selectfont #1}%
\newcommand{\oldB}[1]{\fontspec[Ligatures={Common}]{Didot}\fontsize{12pt}{15pt}\color{solarized@base02}\selectfont #1}%
\newcommand{\tfont}[1]{\fontspec[Alternate=1,Ligatures={Common}]{Hoefler Text}\fontsize{12pt}{20pt}\selectfont #1}%
\newcommand{\dfont}[1]{\fontspec[Ligatures={Common}]{Didot}\fontsize{12pt}{12pt}\selectfont #1}%

\newcommand{\minimize}{\mathop{\mathrm{minimize}}}
\newcommand{\argmin}{\mathop{\mathrm{arg\,min}}}
\newcommand{\argmax}{\mathop{\mathrm{arg\,max}}}
\newcommand{\st}{\mathop{\mathrm{subject\,\,to}}}

\newcommand\independent{\protect\mathpalette{\protect\independenT}{\perp}}
\def\independenT#1#2{\mathrel{\rlap{$#1#2$}\mkern2mu{#1#2}}}

\setlength{\parindent}{0pt}
\setlength{\parskip}{12pt}

\setromanfont [Ligatures={Common}, Numbers={OldStyle}, Variant=01,
 BoldFont={LinLibertine_RB.otf},
 ItalicFont={LinLibertine_RI.otf},
 BoldItalicFont={LinLibertine_RBI.otf}
 ]{LinLibertine_R.otf}



\begin{document}

%%%%%%%%%%%%%%%%%%%%%%%%%%%%%%%%%%%%%%%%%%%%%%%%%%%
\begin{frame}[fragile] \frametitle{}

\vfill

{\fontsize{0.7cm}{0cm}\selectfont Lecture 21 \\\vspace{0.2cm}
Theory of the Lasso II}\\\vspace{0.5cm}
02 December 2015

\vspace{2cm}

\begin{minipage}{0.6\textwidth}
Taylor B. Arnold \\
Yale Statistics  \\
STAT 312/612
\end{minipage}
\hfill
\begin{minipage}{0.3\textwidth}\raggedleft
\includegraphics[scale=0.3]{../yale-logo.png}
\end{minipage}%

\end{frame}

%%%%%%%%%%%%%%%%%%%%%%%%%%%%%%%%%%%%%%%%%%%%%%%%%%%
\begin{frame}[fragile] \frametitle{}

{\color{yaleblue}\fontsize{16pt}{20pt}\selectfont Class Notes}

\begin{itemize}
\item Midterm II - Available now, due next Monday
\item Problem Set 7 - Available now, due December 11th (grace period through the 16th)
\end{itemize}

\end{frame}

%%%%%%%%%%%%%%%%%%%%%%%%%%%%%%%%%%%%%%%%%%%%%%%%%%%
\begin{frame}[fragile] \frametitle{}

\begin{flushright}
{\color{yaleblue}\sc\fontsize{1cm}{0cm}\selectfont Last time}
\end{flushright}

\end{frame}

%%%%%%%%%%%%%%%%%%%%%%%%%%%%%%%%%%%%%%%%%%%%%%%%%%%
\begin{frame}[fragile] \frametitle{}

Last class, we started investigating the theory of the lasso estimator.

For the case of $X^tX$ equal to the identity matrix, we were able to quickly
establish bounds on the prediction error, estimation of $\beta$, and the
reconstruction of the support of $\beta$.

\pause For an arbitrary $X$ matrix we were able to calculate a bound on $|| X (\widehat{\beta} - \beta) ||_2^2$.

\pause Today's goal is to establish a bound on $|| \widehat{\beta} - \beta ||_2^2$


\end{frame}

%%%%%%%%%%%%%%%%%%%%%%%%%%%%%%%%%%%%%%%%%%%%%%%%%%%
\begin{frame}[fragile] \frametitle{}

The basic starting point from last time was the following decomposition, which
had no assumptions beyond linearity of the true model:
 \begin{align*}
\red{|| X (\beta - b) ||_2^2} &\leq \violet{2 \epsilon^t X (b - \beta)} + \blue{\lambda \cdot \left( || \beta ||_1 - || b ||_1 \right)}
\end{align*}
Where can think of this decomposition as \red{the loss to be minimized}, the \violet{empirical part},
and \blue{the penalty term}.

\end{frame}

%%%%%%%%%%%%%%%%%%%%%%%%%%%%%%%%%%%%%%%%%%%%%%%%%%%
\begin{frame}[fragile] \frametitle{}

I then defined the set
\begin{align*}
\mathcal{A} &= \left\{ 2 || \epsilon^t X ||_\infty \leq \lambda \right\}
\end{align*}
And showed that for any $A>1$ we have $\mathbb{P} \mathcal{A} \geq 1 - A^{-1}$
whenever
\begin{align*}
\lambda \geq A \cdot \sqrt{8 \log(2p) \sigma^2}.
\end{align*}

\end{frame}

%%%%%%%%%%%%%%%%%%%%%%%%%%%%%%%%%%%%%%%%%%%%%%%%%%%
\begin{frame}[fragile] \frametitle{}

Today we will motivate a stronger assumption on the model
and use these two results to establish bounds on the prediction of
$\beta$.

Also, it will be helpful to write the set $\mathcal{A}$ as being parameterized
by the value of $\lambda_0$:
\begin{align*}
\mathcal{A}(\lambda_0) &= \left\{ 2 || \epsilon^t X ||_\infty \leq \lambda_0 \right\}
\end{align*}

\end{frame}

%%%%%%%%%%%%%%%%%%%%%%%%%%%%%%%%%%%%%%%%%%%%%%%%%%%
\begin{frame}[fragile] \frametitle{}

\begin{flushright}
{\color{yaleblue}\sc\fontsize{1cm}{0cm}\selectfont Bounds on estimation error}
\end{flushright}

\end{frame}

%%%%%%%%%%%%%%%%%%%%%%%%%%%%%%%%%%%%%%%%%%%%%%%%%%%
\begin{frame}[fragile] \frametitle{}

We already know that on $\mathcal{A}(\lambda_0)$ and with $\lambda > 2 \cdot \lambda_0$, we have:
\begin{align*}
|| X (b - \beta) ||_2^2 + \lambda \cdot || b ||_1 &\leq 2 \epsilon^t X (b - \beta) + \lambda \cdot || \beta ||_1 \\
&\leq \lambda_0 || b - \beta ||_1 + \lambda \cdot || \beta ||_1
\end{align*}
\pause Now, multiplying by two gives:
\begin{align*}
2 || X (b - \beta) ||_2^2 + 2 \lambda \cdot || b ||_1
&\leq \lambda || b - \beta ||_1 + 2 \lambda \cdot || \beta ||_1
\end{align*}

\end{frame}

%%%%%%%%%%%%%%%%%%%%%%%%%%%%%%%%%%%%%%%%%%%%%%%%%%%
\begin{frame}[fragile] \frametitle{}

Recall that we defined the notation: $S = \{j \, : \, \beta_j \neq 0\}$, $s$
is the size of the set $S$, and $v_{S}$ is the vector
$v$ which has components not in $S$ set to zero.

\end{frame}

%%%%%%%%%%%%%%%%%%%%%%%%%%%%%%%%%%%%%%%%%%%%%%%%%%%
\begin{frame}[fragile] \frametitle{}

Notice that:
\begin{align*}
|| b ||_1 &= || b_S ||_1 + || b_{S^c} ||_1 \\
&\geq || \beta ||_1 - || b_S - \beta ||_1 + || b_{S^c} ||_1
\end{align*}
Using the (reverse) triangle inequality and the fact that $\beta_{S^c}$ is zero
by definition.

\end{frame}

%%%%%%%%%%%%%%%%%%%%%%%%%%%%%%%%%%%%%%%%%%%%%%%%%%%
\begin{frame}[fragile] \frametitle{}

Similarly, we have:
\begin{align*}
|| b - \beta ||_1 &= || b_S - \beta_S ||_1 + || b_{S^c} ||_1
\end{align*}
Where clearly $\beta_S$ is redundant, but useful to keep the
notation straight.

\end{frame}

%%%%%%%%%%%%%%%%%%%%%%%%%%%%%%%%%%%%%%%%%%%%%%%%%%%
\begin{frame}[fragile] \frametitle{}

Plugging these in, we now get:
\begin{align*}
2 || X (b - \beta) ||_2^2 + &2 \lambda \cdot || \beta_S ||_1 - 2 \lambda \cdot || b_S - \beta ||_1
  + 2 \lambda \cdot || b_{S^c} ||_1 \\
&\leq \lambda || b - \beta ||_1 + \lambda \cdot || b_S - \beta_S ||_1 + 2 \lambda \cdot || b_{S^c} ||_1
\end{align*}
\pause Which cancels out as:
\begin{align*}
2 || X (b - \beta) ||_2^2 + \lambda || b_{S^c} ||_1 \leq 3 \cdot \lambda \cdot || b_S - \beta_S ||_1
\end{align*}

\end{frame}

%%%%%%%%%%%%%%%%%%%%%%%%%%%%%%%%%%%%%%%%%%%%%%%%%%%
\begin{frame}[fragile] \frametitle{}

This result now actually gives two sub-results, as all three terms are positive and
therefore each component of the left hand side is individually bounded by the right
hand side.

\pause In particular, we have:
\begin{align*}
|| b_{S^c} ||_1 \leq 3 \cdot || b_S - \beta_S ||_1
\end{align*}
Which implies that the amount of error in $b$ can not be too highly concentrated
on $S^c$.

\end{frame}

%%%%%%%%%%%%%%%%%%%%%%%%%%%%%%%%%%%%%%%%%%%%%%%%%%%
\begin{frame}[fragile] \frametitle{}

The other sub-result gives:
\begin{align*}
2 || X (b - \beta) ||_2^2 \leq 3 \lambda \cdot || b_S - \beta_S ||_1
\end{align*}
\pause If $\sigma_{min}$ is the minimum singular value of $X$, then the
left hand side can be bounded below by:
\begin{align*}
2 \sigma^2_{min} || b - \beta ||_2^2 \leq 3 \lambda \cdot || b_S - \beta_S ||_1
\end{align*}
\pause Using the Cauchy-Schwarz inequality, this becomes:
\begin{align*}
2 \sigma^2_{min} || b - \beta ||_2^2 \leq 3 \lambda \cdot \sqrt{s} || b_S - \beta_S ||_2\\
|| b - \beta ||_2 \leq \frac{3 \lambda \sqrt{s}}{2\sigma^2_{min}}
\end{align*}
Which gives a bound on the error of estimating $\beta$, which is exactly what we wanted
to establish.

\end{frame}

%%%%%%%%%%%%%%%%%%%%%%%%%%%%%%%%%%%%%%%%%%%%%%%%%%%
\begin{frame}[fragile] \frametitle{}

Why is this not sufficient for us? Well, in the high dimensional case $p > n$, we will
always have $\sigma_{min}$ equal to $0$.

\pause We can get around this problem by defining a modified version of the minimum
eigenvector (or squared singular value) by only considering $b - \beta$ such that:
\begin{align*}
|| b_{S^c} ||_1 \leq 3 \cdot || b_S - \beta_S ||_1
\end{align*}

\end{frame}

%%%%%%%%%%%%%%%%%%%%%%%%%%%%%%%%%%%%%%%%%%%%%%%%%%%
\begin{frame}[fragile] \frametitle{}

The (minimum) restricted eigenvalue $\phi_S$ on the set $S$ is defined as:
\begin{align*}
\phi_S &= \argmin_{v \, \in \, \mathcal{V}_S} \frac{|| X b ||_2}{|| b ||_2}
\end{align*}
Where:
\begin{align*}
\mathcal{V}_S &= \{ v \in \mathbb{R}^p \, \text{s.t.} \, || v_{S^c} ||_1 \leq 3 \cdot || v_S ||_1  \}
\end{align*}

\end{frame}

%%%%%%%%%%%%%%%%%%%%%%%%%%%%%%%%%%%%%%%%%%%%%%%%%%%
\begin{frame}[fragile] \frametitle{}

Because we do not know $S$, it is impossible to calculate $\phi_S$ in practice.
In theoretical work, often one considers \textbf{the} restricted eigenvalue
$\phi$ defined as the smallest $\phi_S$ for all sets $S$ with size bounded by
some predefined $s_0$.

\end{frame}

%%%%%%%%%%%%%%%%%%%%%%%%%%%%%%%%%%%%%%%%%%%%%%%%%%%
\begin{frame}[fragile] \frametitle{}

Now, we can bound the following using our prior result:
\begin{align*}
2 || X ( b - \beta) ||_2^2 &+ \lambda \cdot || b - \beta ||_1 \\
  &= 2 || X ( b - \beta) ||_2^2 + \lambda \cdot || b_S - \beta_S ||_1 + \lambda \cdot || b_{S^c} ||_1 \\
  &= 4 \lambda \cdot || b_S - \beta_S ||_1
\end{align*}
\pause Using Cauchy-Schartz again, we can change the $\ell_1$-norm to an $\ell_2$-norm at
the cost of a factor of $\sqrt{s}$:
\begin{align*}
2 || X ( b - \beta) ||_2^2 + \lambda \cdot || b - \beta ||_1
  &\leq 4 \lambda \cdot \sqrt{s} \cdot || b_S - \beta_S ||_2
\end{align*}
\pause Finally, we now use the restricted eigenvalue $\phi$ to convert from $\beta$ space
to $X \beta$ space:
\begin{align*}
2 || X ( b - \beta) ||_2^2 + \lambda \cdot || b - \beta ||_1
  &\leq 4 \lambda \cdot \sqrt{s} \cdot || X( b_S - \beta_S) ||_2 / \phi \\
\end{align*}

\end{frame}

%%%%%%%%%%%%%%%%%%%%%%%%%%%%%%%%%%%%%%%%%%%%%%%%%%%
\begin{frame}[fragile] \frametitle{}

I am now going to use an inequality trick that is often useful in
theoretical statistics derivations. For any $u$ and $v$, notice that
$4uv \leq u^2 + 4 v^2$.

For a proof, notice that it is trivially true at zero and negative values
of $u$ and $v$. Then look at the derivatives and notice that the right hand
side grows faster than the left hand side in the directions of both $u$ and
$v$.

\end{frame}


%%%%%%%%%%%%%%%%%%%%%%%%%%%%%%%%%%%%%%%%%%%%%%%%%%%
\begin{frame}[fragile] \frametitle{}

Setting $u = || X( b_S - \beta_S) ||_2$, we then have:
\begin{align*}
2 || X ( b - \beta) ||_2^2 + \lambda \cdot || b - \beta ||_1
  &\leq || X( b_S - \beta_S) ||_2 + 4 \lambda^2 \cdot s \cdot / \phi^2 \\
  &\leq || X( b - \beta) ||_2^2 + 4 \lambda^2 \cdot s \cdot / \phi^2
\end{align*}
\pause And when canceling one factor of $|| X( b - \beta) ||_2$:
\begin{align*}
|| X ( b - \beta) ||_2^2 + \lambda \cdot || b - \beta ||_1
  &\leq 4 \lambda^2 \cdot s \cdot / \phi^2
\end{align*}
Which holds on the entire set $\mathcal{A}(\lambda_0)$.

\end{frame}

%%%%%%%%%%%%%%%%%%%%%%%%%%%%%%%%%%%%%%%%%%%%%%%%%%%
\begin{frame}[fragile] \frametitle{}

This establishes two simultaneous bounds:
\begin{align*}
|| X ( b - \beta) ||_2^2 &\leq 4 \lambda^2 \cdot s \cdot / \phi^2 \\
|| b - \beta ||_1 &\leq 4 \lambda \cdot s \cdot / \phi^2
\end{align*}
Though the first is slightly less satisfying than our result
in last class as it relies on $\phi^2$, though it no longer requires
the norm of $\beta$.

\end{frame}

%%%%%%%%%%%%%%%%%%%%%%%%%%%%%%%%%%%%%%%%%%%%%%%%%%%
\begin{frame}[fragile] \frametitle{}

\begin{flushright}
{\color{yaleblue}\sc\fontsize{1cm}{0cm}\selectfont Asymptotic analysis}
\end{flushright}

\end{frame}

%%%%%%%%%%%%%%%%%%%%%%%%%%%%%%%%%%%%%%%%%%%%%%%%%%%
\begin{frame}[fragile] \frametitle{}

As before, we can convert a more natural re-scaled problem by dividing all of the
$\lambda$ parameters by $\sqrt{n}$

Also, remember that for some $A > 1$, we have $\mathbb{P} \mathcal{A}(\lambda_0) \geq 1 - A^{-1}$
for all $\lambda > A \cdot \sqrt{16 n^{-1} \log(2p) \sigma^2}$.

\end{frame}

%%%%%%%%%%%%%%%%%%%%%%%%%%%%%%%%%%%%%%%%%%%%%%%%%%%
\begin{frame}[fragile] \frametitle{}

Therefore, we have:
\begin{align*}
|| b - \beta ||_1 &\leq 4 \lambda \cdot s \cdot / \phi^2 \\
&\leq 8 \cdot A \sigma^2 / \phi^2 \cdot \frac{s_n^2 \log(2p_n)}{n}
\end{align*}
Which is the same result as from the Bickel, Ritov, Tsybakov paper.

\end{frame}

%%%%%%%%%%%%%%%%%%%%%%%%%%%%%%%%%%%%%%%%%%%%%%%%%%%
\begin{frame}[fragile] \frametitle{}

To establish consistency of the estimator under constant noise and
restricted eigenvalues $\phi^2$, we need the following limit to
go to zero:
\begin{align*}
\lim_{n \rightarrow \infty} \frac{s_n^2 \log(2p_n)}{n} = 0
\end{align*}
Which can happen with a number of different scalings, such as a constant
number of non-zero terms but an exponential number of non-zero terms. Or,
$s_n$ growing like $n^{1/3}$ and $p_n$ growing linearly with $s_n$.

\end{frame}


%%%%%%%%%%%%%%%%%%%%%%%%%%%%%%%%%%%%%%%%%%%%%%%%%%%
\begin{frame}[fragile] \frametitle{}

Some (personal) closing thoughts on the application of the lasso theory
to data analysis:
\begin{enumerate}
\item The theory is useful for establishing a rough rule of thumb
for how large $p_n$ and $s_n$ can be to have a reasonable chance
of reconstructing $\beta$ or $X\beta$ \pause
\item The theory also helps guide where to start looking for the
optimal $\lambda$ \pause
\item We still generally need some form of cross validation however,
as the theoretical values tend to overestimate $\lambda$ in
practice; we also do not know $\sigma^2$ and in theory need to use
an over-estimate for the convergence results to hold \pause
\item Bounds on $|| X (\beta - \widehat{\beta}) ||_2^2$ are nice to have,
however the theoretical bounds on $|| \beta - \widehat{\beta} ||_2^2$
are difficult to use in practice due to the near-impossible to calculate
restricted eigenvalue assumption \pause
\item I have always been skeptical of the asymptotic results for the same
reason; $\phi$ likely depends on $n$, $p_n$ and $s_n$ in complex ways that
are not accounted for
\end{enumerate}

\end{frame}

%%%%%%%%%%%%%%%%%%%%%%%%%%%%%%%%%%%%%%%%%%%%%%%%%%%
\begin{frame}[fragile] \frametitle{}

For our next (and last) week we will:
\begin{enumerate}
\item use the lasso to encode more complex forms of linear sparsity (e.g., outlier
detection and the fused lasso)
\item give an alternative approach to solving for the lasso solution at
a particular value of $\lambda$
\end{enumerate}

\end{frame}

\end{document}











